%% kapa: exam=1 round=28 year=2017 period=1 form=A subject=economics subject_ko=경제학원론 seq=26 number=66
%% source: Q-Net 감정평가사 기출문제 (cst00309) artlSeq=5199603 file=extracted/1008/5199603/2017년 제28회 감정평가사 1차시험 문제지/1교시 A형.pdf
%% conversion: kordoc 4.24.51 -> tools/gen1.py (auto); review=proofread-sonnet; flags=none
\begin{question}{66}
통화량 변동에 관한 설명으로 옳지 않은 것은?
\begin{choices}
  \item 법정지급준비율의 변동은 본원통화량을 변화시키지 않는다.
  \item 중앙은행이 통화안정증권을 발행하여 시장에 매각하면 통화량이 감소한다.
  \item 중앙은행이 시중은행으로부터 채권을 매입하면 통화량이 감소한다.
  \item 은행의 법정지급준비율을 100 \%로 규제한다면 본원통화량과 통화량은 동일하다.
  \item 정부의 중앙은행차입이 증가하면 통화량은 증가한다.
\end{choices}
\end{question}
\answer{3}
